\documentclass[10pt,a4paper]{amsart}
\usepackage{amsmath,amssymb,amsthm,graphicx}
% Retained minimal preamble for the Stacks Project chapter more-algebra.tex (More on Algebra).
% The chapter file pinned-original-source.tex carries only \input{preamble} on its first line, so no
% preamble text lives inside the chapter itself.  The Stacks Project uses one single shared preamble for
% every chapter; the snapshot stages it identically as sources/stacks-chow/preamble.tex and
% sources/stacks-etale-cohomology/preamble.tex.  The block below is copied mechanically, line for line,
% from sources/stacks-chow/preamble.tex lines 14-18 (the Xy-pic setup) and 154-217 (the native theorem
% declarations together with the mathematical macros).  Nothing was added, removed or redefined; the
% wrapper class and the amsmath/amssymb/amsthm/graphicx base packages are supplied by the assembler.
% Chapter: more-algebra.tex -- More on Algebra.
% For commutative diagrams we use Xy-pic
\usepackage[all]{xy}

% We use 2cell for 2-commutative diagrams.
\xyoption{2cell}

% Theorem environments.
%
\theoremstyle{plain}
\newtheorem{theorem}[subsection]{Theorem}
\newtheorem{proposition}[subsection]{Proposition}
\newtheorem{lemma}[subsection]{Lemma}

\theoremstyle{definition}
\newtheorem{definition}[subsection]{Definition}
\newtheorem{example}[subsection]{Example}
\newtheorem{exercise}[subsection]{Exercise}
\newtheorem{situation}[subsection]{Situation}

\theoremstyle{remark}
\newtheorem{remark}[subsection]{Remark}
\newtheorem{remarks}[subsection]{Remarks}

\numberwithin{equation}{subsection}

% Macros
%
\def\lim{\mathop{\mathrm{lim}}\nolimits}
\def\colim{\mathop{\mathrm{colim}}\nolimits}
\def\Spec{\mathop{\mathrm{Spec}}}
\def\Hom{\mathop{\mathrm{Hom}}\nolimits}
\def\Ext{\mathop{\mathrm{Ext}}\nolimits}
\def\SheafHom{\mathop{\mathcal{H}\!\mathit{om}}\nolimits}
\def\SheafExt{\mathop{\mathcal{E}\!\mathit{xt}}\nolimits}
\def\Sch{\mathit{Sch}}
\def\Mor{\mathop{\mathrm{Mor}}\nolimits}
\def\Ob{\mathop{\mathrm{Ob}}\nolimits}
\def\Sh{\mathop{\mathit{Sh}}\nolimits}
\def\NL{\mathop{N\!L}\nolimits}
\def\CH{\mathop{\mathrm{CH}}\nolimits}
\def\proetale{{pro\text{-}\acute{e}tale}}
\def\etale{{\acute{e}tale}}
\def\QCoh{\mathit{QCoh}}
\def\Ker{\mathop{\mathrm{Ker}}}
\def\Im{\mathop{\mathrm{Im}}}
\def\Coker{\mathop{\mathrm{Coker}}}
\def\Coim{\mathop{\mathrm{Coim}}}

% Boxtimes
%
\DeclareMathSymbol{\boxtimes}{\mathbin}{AMSa}{"02}

%
% Macros for moduli stacks/spaces
%
\def\QCohstack{\mathcal{QC}\!\mathit{oh}}
\def\Cohstack{\mathcal{C}\!\mathit{oh}}
\def\Spacesstack{\mathcal{S}\!\mathit{paces}}
\def\Quotfunctor{\mathrm{Quot}}
\def\Hilbfunctor{\mathrm{Hilb}}
\def\Curvesstack{\mathcal{C}\!\mathit{urves}}
\def\Polarizedstack{\mathcal{P}\!\mathit{olarized}}
\def\Complexesstack{\mathcal{C}\!\mathit{omplexes}}
% \Pic is the operator that assigns to X its picard group, usage \Pic(X)
% \Picardstack_{X/B} denotes the Picard stack of X over B
% \Picardfunctor_{X/B} denotes the Picard functor of X over B
\def\Pic{\mathop{\mathrm{Pic}}\nolimits}
\def\Picardstack{\mathcal{P}\!\mathit{ic}}
\def\Picardfunctor{\mathrm{Pic}}
\def\Deformationcategory{\mathcal{D}\!\mathit{ef}}
\usepackage{cleveref}
\crefname{originallemma}{original lemma}{original lemmas}
\Crefname{originallemma}{Original lemma}{Original lemmas}
\makeatletter
\newlabel{lemma-finite-gl-dim-tor-dimension}{{066P}{source}{}{sourcepointer-1}{}}
\newlabel{lemma-finite-gl-dim-tor-dimension@cref}{{[originallemma][0][]066P}{[1][0][]source}{}{}{}}
\makeatother
\crefname{originallemma}{original lemma}{original lemmas}
\Crefname{originallemma}{Original lemma}{Original lemmas}
\makeatletter
\newlabel{lemma-check-perfect-pointwise}{{068V}{source}{}{sourcepointer-2}{}}
\newlabel{lemma-check-perfect-pointwise@cref}{{[originallemma][0][]068V}{[1][0][]source}{}{}{}}
\makeatother
\title{Pseudo-coherent complexes with bounded tor amplitude over a fibrewise finite-global-dimension base are perfect}
\author{The Stacks Project Authors}
\date{stacks-project:more-algebra.tex:local-snapshot-20261012; GNU FDL 1.2 or later}
\begin{document}
\maketitle
\setcounter{subsection}{1}
\noindent\textit{Editorial note.} All delivered mathematical text is the author's own, copied from the pinned Stacks Project chapter file more-algebra.tex; no statement or proof step was rewritten, reordered or removed. The counted statement is the labeled lemma environment at source lines 20174-20188 (canonical tag 068X) and the counted proof is the proof environment at source lines 20190-20221, which immediately follows it and is its own proof.\par\medskip
\noindent\textit{Editorial note.} The proof refers twice inside the chapter: to lemma-finite-gl-dim-tor-dimension at source line 17287 for the fibrewise tor amplitude estimate and to lemma-check-perfect-pointwise at source line 20032 for the fibrewise criterion of perfectness. Both are navigation pointers and their roles are recorded in the external reference reasons; the displayed computation of the fibres over primes of B is inside the counted span.\par\medskip
\noindent\textit{Editorial note.} The lead-in at source lines 20170-20172, announcing the case of a polynomial ring extension, is outside the counted spans, as is the section text that follows the proof.\par\medskip
\noindent\textit{Editorial note.} The chapter file pinned-original-source.tex carries only an input of the shared preamble on its first line, so no preamble text lives inside the chapter. The retained file sources/stacks-more-algebra/retained-preamble.tex is a mechanical, line-for-line copy of the project's single shared preamble as staged in the snapshot (Xy-pic setup plus the native theorem declarations and mathematical macros). Nothing was added, removed or redefined, and the excerpt serves only as preamble, never as a span coordinate.\par\medskip
\noindent\textit{Editorial note.} The source URL is the source-level Stacks Project URL, as required when more than one card is produced from one source; the per-claim tag is carried by claim\_key and by this note. The counted statement and proof cite no bibliography entry, so this card needs no bibliography.\par\medskip
\noindent\textit{Allowed standard prerequisites (editorial).}\begin{itemize}
\item A complex over a ring of global dimension at most d whose cohomology is concentrated in degrees a to b has tor amplitude in the shifted range a - d to b. invoked at proof lines 20201-20204, where the fibre K tensor\_A kappa(p) over the residue field is estimated using the finite global dimension hypothesis. Converts the fibrewise cohomology bound into the fibrewise tor amplitude bound that has to be propagated to the fibres over primes of B.
\item Perfectness of a pseudo-coherent complex over a ring can be checked on the fibres at the primes, by testing the vanishing of the derived tensor products with the residue fields. invoked at proof lines 20217-20220, where the vanishing of the shifted cohomology for all primes q of B gives the conclusion. Turns the pointwise computation just performed into the global perfectness statement with the stated tor amplitude range.
\item For a flat ring map A to B and a complex of B-modules that is flat over A, the derived tensor product with a B-module agrees with the ordinary tensor product, and localization at a prime of B is computed through the residue field of the prime below. used at proof lines 20191-20205 and 20206-20216. Justifies replacing the derived tensors over A and over B by ordinary ones and the displayed chain of identifications for K tensor\textasciicircum{}L\_B kappa(q).
\end{itemize}\medskip
\ifdefined\hypertarget\hypertarget{sourcepointer-1}{}\fi
\noindent\textit{Source reference.} Original lemma 066P: sources/stacks-more-algebra/pinned-original-source.tex, source lines 17287--17287. The reference is used at proof lines 20201-20204 to turn the fibrewise cohomology bound into the fibrewise tor amplitude bound over A; the rest of the proof, namely the reduction to finite free complexes and the computation of the fibres over B, is inside the counted span.\par\medskip
\ifdefined\hypertarget\hypertarget{sourcepointer-2}{}\fi
\noindent\textit{Source reference.} Original lemma 068V: sources/stacks-more-algebra/pinned-original-source.tex, source lines 20032--20032. The reference is the fibrewise criterion invoked at proof lines 20217-20220 to conclude perfectness; the vanishing of the cohomology of every fibre, which is the condition the criterion tests, is established inside the counted span at lines 20205-20216.\par\medskip
\begin{lemma}
\label{lemma-perfect-over-polynomial-ring}
Let $A \to B$ be a ring map. Let $a, b \in \mathbf{Z}$. Let $d \geq 0$.
Let $K^\bullet$ be a complex of $B$-modules. Assume
\begin{enumerate}
\item the ring map $A \to B$ is flat,
\item for every prime $\mathfrak p \subset A$ the ring
$B \otimes_A \kappa(\mathfrak p)$ has finite global dimension $\leq d$,
\item $K^\bullet$ is pseudo-coherent as a complex of $B$-modules, and
\item $K^\bullet$ has tor amplitude in $[a, b]$ as a complex
of $A$-modules.
\end{enumerate}
Then $K^\bullet$ is perfect as a complex of $B$-modules
with tor amplitude in $[a - d, b]$.
\end{lemma}

\begin{proof}
We may assume that $K^\bullet$ is a bounded above complex of
finite free $B$-modules. In particular, $K^\bullet$ is flat as a
complex of $A$-modules and
$K^\bullet \otimes_A M = K^\bullet \otimes_A^{\mathbf{L}} M$ for any
$A$-module $M$. For every prime $\mathfrak p$ of $A$ the complex
$$
K^\bullet \otimes_A \kappa(\mathfrak p)
$$
is a bounded above complex of finite free modules over
$B \otimes_A \kappa(\mathfrak p)$ with vanishing $H^i$ except
for $i \in [a, b]$. As $B \otimes_A \kappa(\mathfrak p)$
has global dimension $d$ we see from
Lemma \ref{lemma-finite-gl-dim-tor-dimension}
that $K^\bullet \otimes_A \kappa(\mathfrak p)$ has tor amplitude in
$[a - d, b]$. Let $\mathfrak q$ be a prime of $B$ lying over $\mathfrak p$.
Since $K^\bullet \otimes_A \kappa(\mathfrak p)$ is a bounded above
complex of free $B \otimes_A \kappa(\mathfrak p)$-modules we see
that
\begin{align*}
K^\bullet \otimes_B^{\mathbf{L}} \kappa(\mathfrak q)
& = K^\bullet \otimes_B \kappa(\mathfrak q) \\
& = (K^\bullet \otimes_A \kappa(\mathfrak p))
\otimes_{B \otimes_A \kappa(\mathfrak p)} \kappa(\mathfrak q) \\
& = (K^\bullet \otimes_A \kappa(\mathfrak p))
\otimes^{\mathbf{L}}_{B \otimes_A \kappa(\mathfrak p)} \kappa(\mathfrak q)
\end{align*}
Hence the arguments above imply that
$H^i(K^\bullet \otimes_B^{\mathbf{L}} \kappa(\mathfrak q)) = 0$
for $i \not \in [a - d, b]$. We conclude by
Lemma \ref{lemma-check-perfect-pointwise}.
\end{proof}

\end{document}
